% ==========================================================================
% GENERATED by scripts/split_paper.py. Do not edit this file.
%
% Source: imported_papers/Neural_Computation__CNA/NECO-D-24-00304-Source.tex
% Part 10: Feature Extraction
%
% Edits belong upstream, in the published paper. To change what the
% chapter includes, change imports/cna.toml and re-run
% the splitter.
% ==========================================================================
\section{Feature Extraction}\label{sec:features}
\begin{figure}[ht!]
  \begin{center}
    \includegraphics[width=0.7\textwidth]{NECO-D-24-00304-Figure.9}
  \end{center}
  \caption{\textbf{Features}: The hand-crafted convolutional kernels of the feature extractor.}\label{fig:weights_feature_extractor}
\end{figure}

The stage $S1$ serves to extract initial binary features from the images.
The input images have $C=1$ channel and $C^{(S1)}=4$ features are extracted at each spatial location. We utilize hand-crafted filters with a size of $5 \times 5$ as displayed in Figure \ref{fig:weights_feature_extractor}, each of the four filters corresponding to a particular line orientation (vertical, positive diagonal, horizontal, and negative diagonal).
This approach is motivated by the observation that straight lines when viewed locally, can be expressed as combinations of these fundamental line types (e.g., a line with a $20^{\circ}$ angle activates the horizontal and the $+45^{\circ}$ diagonal feature neurons). The use of such manually designed filters allows a more straightforward interpretation of the results. While more complex datasets necessitate learned filters, our primary focus in this work is on constructing net fragments, and thus, we favor the simplicity of hand-crafted filters.

These filters move across the image with a step size (`stride') of $1$ (assuming $0$ as input signal beyond the image border: `zero padding') so that the input image and output image are of the same size.
Thus, the input $\boldsymbol{x}$ in the feature extraction stage are images of size $(1\times32\times32)$ and the output $\boldsymbol{y}^{(S1)}$ feature activation maps of size $(4\times32\times32)$. The dataset's lines activate multiple filters at different positions, facilitating the construction of net fragments across these feature channels.
We convert the floating-point output of the convolutional weights to binary activations using the heaviside function $\Theta$ with a bias of $b^{(S1)}=0.5$, i.e. neurons fire a binary spike if $\boldsymbol{a}_j > 0.5$ (see eq. \ref{eq:y_s1}).
